\section{Zahnradschäden und Zahnfußtragfähigkeit}
\label{sec:schaeden}

VDI 2736 Blatt 1 ordnet die Versagensarten von Kunststoffzahnrädern den Einsatzbedingungen zu. Im Trockenlauf kommt es oberhalb einer werkstoffabhängigen Last zu Anschmelzungen, die innerhalb kurzer Zeit zur völligen Zerstörung der Zahnflanken führen. Daneben sind Zahnbrüche und Verschleiß möglich \cite[S.~14]{VDI2736-1}. Zahnbrüche können sowohl von der aktiven Flanke als auch von außerhalb ausgehen, Verschleiß verändert die Zahnoberfläche und verringert zugleich den Zahnquerschnitt. Im Betrieb mit externer Schmierung treten vorwiegend Zahnbrüche und Grübchenbildung auf \cite[S.~15]{VDI2736-1}. Die Tabelle der Richtlinie führt Anschmelzungen durch zu starke Erwärmung der Zahnflanken und den Zahnfußbruch durch zu hohe Biegebeanspruchung auf. Zahnflankenbruch und Grübchenbildung entstehen durch zu hohe Hertzsche Beanspruchung mit überlagerter Schub- und Biegebeanspruchung. Zahnverschleiß folgt aus zu geringer Verschleißfestigkeit oder der abrasiven Wirkung des Reibpartners, Deformation aus einer Beanspruchung über der Verformungsgrenze \cite[S.~16]{VDI2736-1}. Niemann und Winter halten fest, dass an Kunststoffzahnrädern dieselben Schäden wie an Stahlzahnrädern auftreten können, nämlich Grübchen und Verschleiß sowie Zahnbruch \cite[S.~350]{NieWin03}. Die Tragfähigkeitsberechnung nach VDI 2736 Blatt 2 lehnt sich an DIN 3990 an und berücksichtigt die Besonderheiten der Thermoplaste durch stoffspezifische Faktoren \cite[S.~5 f.]{VDI2736-2}. Dazu zählen die Abhängigkeit der Festigkeit von der Betriebstemperatur sowie von Höhe und Dauer der Belastung und von der Belastungsgeschwindigkeit, die schlechtere Wärmeleitfähigkeit und die größere Deformation. Die Nachrechnung für Ermüdungsbeanspruchung umfasst die Zahntemperatur, die Zahnfußtragfähigkeit, die Zahnflankentragfähigkeit gegen Grübchen und Verschleiß sowie die Verformung, ergänzt um eine Nachrechnung für Spitzenlasten \cite[S.~6]{VDI2736-2}.

Der Zahnfußbruch führt in der Regel zum sofortigen Totalausfall des Getriebes, weshalb die Lebensdauer von Kunststoffzahnrädern maßgeblich durch die Zahnfußtragfähigkeit bestimmt wird \cite[S.~8]{Ill22}, \cite[S.~5]{Has18}. Nach Linke richtet sich der Nachweis der Zahnfußtragfähigkeit sowohl auf den Ermüdungsbruch als auch auf bleibende Verformung, Anriss oder Gewaltbruch bei maximaler Belastung \cite[S.~297]{Lin10}. Zur Berechnung wird der Zahn als einseitig eingespannter Balken unter reiner Biegung betrachtet, Schub- und Druckanteile sind vernachlässigbar \cite[S.~8]{Ill22}. Im Allgemeinen geht der Dauerbruch von der Zugseite aus, etwa an der Stelle der Berührung mit der $30$°-Tangente \cite[S.~147]{NieWin03}. Die maximale örtliche Zugspannung an dieser Stelle gilt deshalb als Maß für die Bruchgefährdung. Bei Überschreiten der lokalen Biegeschwellfestigkeit entstehen dort Anrisse. Der Riss wächst ins Werkstoffinnere, bis ein Restgewaltbruch den Zahn abtrennt \cite[S.~8]{Ill22}. Als Bezugsquerschnitt gilt bei genauen Analysen der Querschnitt, der sich durch die $30$°-Tangente an die Fußkurve der höchstbelasteten Flanke ergibt \cite[S.~167]{Lin10}. ISO 6336-3 definiert die Sehne zwischen den Berührpunkten der $30$°-Tangenten an die Fußrundungen als kritischen Querschnitt der Berechnung \cite[S.~14]{ISO6336-3}. Auch Kassem et al. nennen den Querschnitt der $30$°-Tangente als Ort des Maximums der Zahnfußspannung, beeinflusst von der Kerbwirkung der Fußrundung \cite[S.~2]{Kas22}. Eine Auswertung von 128 Brüchen an Rädern aus POM und PA46 deutet nach Leipersberger auf einen Rissausgang nahe einer $42$°-Tangente hin \cite[S.~8]{Fie23}. Hasl erkennt darin keinen Einfluss von Baugröße, Last oder Temperatur auf den Tangentenwinkel.

In den Laufversuchen von Hasl entwickelt sich der Zahnfußbruch der Prüfräder aus POM aus einem Anriss in der zugseitigen Fußrundung \cite[S.~93]{Has18}. Rastlinien sind auf den Bruchflächen nicht erkennbar, was für eine kurzfristige Entwicklung des Schadens spricht. Im Bereich hoher Drehmomente geht die Schadensart in ein Versagen der Zahnflanke im fußnahen Bereich über, bei der Verzahnung kst-A überlagert von Anschmelzungen \cite[S.~93 f.]{Has18}. Fiederling beobachtet an der Verzahnung kst-E aus POM bei Drehmomenten bis $12$\,Nm vollständige Zahnfußbrüche, bei $14$\,Nm und mehr dagegen Anrisse mit stetigem Risswachstum und stetig abfallendem Drehmoment (\figref{fig:quelle_fie23_zahnfussbruch}) \cite[S.~66 f.]{Fie23}. Die ersten feinen Risse entstehen rechts und links der Zahnbreitenmitte und vereinigen sich zu einem mittigen Riss, was auf eine Spannungsüberhöhung durch die Radkörpergestaltung des Stegrads zurückgeführt wird \cite[S.~67]{Fie23}. Prüfräder aus PEEK fallen dagegen nicht durch Zahnfußbruch aus, sondern durch grübchenartige Schäden der Flanke und ein Abscheren der Zahnflanke im Bereich des Fußnutzkreises \cite[S.~135]{Has18}, \cite[S.~76]{Fie23}. An PEEK-GF30 treten vornehmlich Grübchen und Verschleiß auf, wobei der Verschleiß die Enden der Glasfasern freilegt \cite[S.~78]{Fie23}.

\begin{figure}[htbp]
    \centering
    \includegraphics[width=0.8\textwidth]{kap2/quelle_fie23_zahnfussbruch}
    \caption{Zahnfußbruch an der Verzahnung kst-E aus POM im Laufversuch: vollständiger Bruch bei $8$\,Nm nach etwa 39,4 Millionen Lastwechseln (links) und Anriss bei $18$\,Nm nach etwa 787\,000 Lastwechseln (rechts), nach \cite[S.~66]{Fie23}.}
    \label{fig:quelle_fie23_zahnfussbruch}
\end{figure}

Grübchen entstehen durch Werkstoffermüdung im oberflächennahen Bereich der aktiven Flanke, wenn die lokale Wälzfestigkeit überschritten wird \cite[S.~9]{Ill22}. In feine Haarrisse wird der Schmierstoff beim Überwälzen eingepresst und sprengt Werkstoffpartikel aus der Oberfläche. Progressives Grübchenwachstum zerrüttet die Flanke und kann über einen Bruch der Zahnflanke zum Totalausfall führen. Abrasiver Verschleiß tritt vor allem im Trockenlauf auf. Starker Verschleiß verringert den Zahnquerschnitt bis zum Bruch des verbleibenden Querschnitts und vergrößert das Verdrehflankenspiel. Unter Ölschmierung wird er durch die Trennung der Oberflächen wirksam verringert \cite[S.~9 f.]{Ill22}. Durch Reibung im Zahnkontakt und durch die viskoelastischen Verluste erwärmen sich Kunststoffzahnräder im Betrieb. Bei Überschreiten der Schmelztemperatur kommt es zum Anschmelzen, das die Zahnflanken zerstört und auch unter geschmierten Bedingungen bei hoher Antriebsleistung auftreten kann \cite[S.~10]{Ill22}. Fürstenberger beobachtet im Trockenlauf, dass der Ausfall des Kunststoffrades innerhalb weniger Minuten nach Beginn der Schadensentstehung erfolgt und bei Erreichen der thermischen Leistungsgrenze mit Anschmelzungen der Zahnflanke einhergeht \cite[S.~73 und 125]{Fuer13}. Große Zahnverformungen wirken sich nach VDI 2736 Blatt 2 wie Teilungsabweichungen aus und verursachen Eingriffsstörungen. Als Orientierungswert für die zulässige Verformung gilt $0{,}07\,m_{\mathrm{n}}$ \cite[S.~19 f.]{VDI2736-2}. Nach Winkler, zitiert bei Hasl \cite[S.~16]{Has18}, geht bei fettgeschmierten Rädern aus PA46, PA66 sowie POM die Schadensart zwischen etwa $10^{3}$ und $4\cdot10^{4}$ Lastspielen von plastischer Verformung zu Zahnfußbruch über.

% TK erledigt (als Sekundärzitat gekennzeichnet, Erstquelle als Platzhalter): Mir fällt auf, dass du manchmal wie im letzten Satz Winkler beobachtet schreibst, aber die Quelle beispielsweise die Arbeit von Hasl ist in der er dann Winkler zitiert. Aus der Quelle Winkler, A. “Experimentelle und theoretische Untersuchungen zur kunststoffgerechten Auslegung von Zahnrädern”. PhD thesis. TU Chemnitz, 2009. Das ist dann kein korrekter Verweis und sollte der unabhägige review dann gegenprüfen, wenn ich fertig bin. Es soll immer die Erstquelle zitiert werden, weil die Aussage in Hasl ein Zitat bezugnehmend auf die Arbeit von Winkler ist. Es gibt sicer noch mehr solche Stellen in der Arbeit. Bitte aufmerksam schauen, falls ich mal eine vergesse anzustreichen.

Für die Ermittlung von Tragfähigkeitskennwerten zählt VDI 2736 Blatt 4 die Lastwechsel bis zum Bruch des Zahns beziehungsweise bis zum Auftreten einer bestimmten Menge an Grübchen. Welche Schadensart für einen Versuch gilt, kann eventuell erst nach Abschluss des Prüflaufs entschieden werden \cite[S.~6]{VDI2736-4}. Als Versagenskriterium für Grübchen gilt ein Flächenanteil von $30$\,\% auf mindestens einer genutzten Zahnflanke \cite[S.~8]{VDI2736-4}. Am nicht mechanisch verspannten Prüfstand gilt das Zahnrad als ausgefallen, wenn die Drehzahl abtriebsseitig entsprechend stark abfällt. Zwei Messzellen erkennen einen Ausfall bereits bei einer nicht mehr exakten Übersetzung \cite[S.~10]{VDI2736-4}. Die Ausfallkriterien von Hasl und Fiederling am Verspannungsprüfstand stehen in \secref{sec:experiment_stand}. Welche Versagensdefinition diese Arbeit verwendet und wie sie auf die Lebensdauer wirkt, wird in \secref{sec:auswertung_laufversuche} festgelegt.

Die Zahnfußtemperatur kann die Zahnfußtragfähigkeit herabsetzen. Die Zeitschwellfestigkeit nach VDI 2736 Blatt 2 wird deshalb für die Zahnfußtemperatur angegeben und fällt mit steigender Temperatur \cite[S.~8 und 14 f.]{VDI2736-2}. Hasl prüft alle Werkstoffe bei einer Zahnfußtemperatur von etwa $80$\,\textdegree C und findet für PA46 eine Zahnfußfestigkeit deutlich unter der von POM \cite[S.~91 und 134]{Has18}. Bei $120$\,\textdegree C liegt die Zahnfußfestigkeit von PA46 nur moderat unter der bei $80$\,\textdegree C. Im Bereich nahe der rechnerischen Hauptresonanz ist die Zahnfußbeanspruchbarkeit der Räder aus POM signifikant geringer als im unter- oder überkritischen Bereich \cite[S.~134]{Has18}. Ölschmierung verbessert die Wärmeabfuhr aus dem Zahnkontakt, so dass bei guten Schmier- und Kühlbedingungen deutlich höhere Drehzahlen und Leistungen möglich sind als im Trockenlauf \cite[S.~10 und 24]{Ill22}, \cite[S.~135]{Has18}. Faserverstärkte Werkstoffe zeigen nach Illenberger vor allem bei hohen Lasten und geringen Laufzeiten Vorteile. Bei hohen Lastspielzahlen und niedrigen Lasten kann ein unverstärkter Werkstoff dagegen eine höhere Lebensdauer erreichen \cite[S.~22]{Ill22}. Weil die Fasern die Biegewechselfestigkeit steigern, kohlenstofffaserverstärktes PEEK aber eine geringere Grübchentragfähigkeit und abrasiven Verschleiß zeigt, schlägt Illenberger Hybridräder mit faserverstärktem Zahnfuß und Zahnkern und unverstärkter Flankenoberfläche vor \cite[S.~106]{Ill22}. Der Steifigkeitsunterschied zwischen der hochorientierten Randschicht und dem Zahninneren erhöht nach FVA 856 I die Zahnfußspannung gegenüber einem isotropen Modell (\secref{sec:spritzgiessen}).

Von der Verzahnungsgeometrie wirken vor allem der Modul, der Eingriffswinkel, die Zähnezahl und die Fußrundung auf die Zahnfußtragfähigkeit. Durch die Wahl des Moduls wird maßgeblich der Zahnfußquerschnitt eines Zahns bestimmt, weshalb die Zahnfußtragfähigkeit das entscheidende Kriterium für den Mindestmodul ist \cite[S.~93]{KloBre17}. Eine Erhöhung des Eingriffswinkels senkt durch die größeren Krümmungsradien und Zahnfußdickensehnen die Beanspruchung bei gleicher Belastung. Zu kleine Eingriffswinkel führen zu Unterschnitt im Zahnfuß und begünstigen den Zahnfußbruch \cite[S.~31 und 95]{KloBre17}. Die Zahnfußdicke an der $30$°-Tangente hängt zudem vom Erzeugungswerkzeug ab \cite[S.~121]{Lin10}. Bei Kunststoffzahnrädern kommt die Überdeckung unter Last hinzu, denn die ertragbaren Drehmomente der Prüfverzahnungen mit größeren Zähnezahlen liegen bei Hasl signifikant über der Erwartung nach VDI 2736 \cite[S.~134]{Has18}. Die Kerbwirkung der Fußrundung bestimmt die Spannungsüberhöhung im Zahnfuß \cite[S.~2]{Kas22}. Fußgeometrien, die diese Kerbwirkung verringern, werden in \secref{sec:numerik_fussgeometrien} behandelt. Verzahnungsabweichungen beeinflussen das Laufverhalten, mit wachsender Teilungs-, Profilwinkel- und Rundlaufabweichung steigt die Verschleißneigung, und größere Zahnformabweichungen verstärken die Geräuschentwicklung \cite[S.~30]{VDI2736-1}.

% Herkunft:
% Absatz 1 (Schadensarten): VDI 2736 Blatt 1 S. 14 (6.2 Versagensarten: Anschmelzungen im Trockenlauf,
%   Zahnbrüche, Verschleiß), S. 15 (Bruchausgang, Verschleiß, externe Schmierung: Zahnbrüche und Grübchen),
%   S. 16 (Tabelle 4 Schadensart und Ursache); Niemann/Winter 2003 S. 350 (22.4 Kunststoffzahnräder: gleiche
%   Schäden wie Stahl: Grübchen, Verschleiß, Zahnbruch); VDI 2736 Blatt 2 S. 5 f. (Abschnitt 3: DIN 3990,
%   Besonderheiten der Thermoplaste, Nachrechnung 3.2.1 bis 3.2.5 und 3.3).
% Absatz 2 (Zahnfußbruch, 30°-Tangente): Illenberger 2022 S. 8 (2.2.1: Totalausfall, Balkenmodell, Schub und
%   Druck vernachlässigbar nach Krause, 30°-Tangente, Anriss, Restgewaltbruch, Bild 2-1); Hasl 2018 S. 5
%   (sofortiger Ausfall); Linke 2010 S. 297 (6.5.2: Ermüdungsbruch, bleibende Verformung, Anriss, Gewaltbruch;
%   Textseite im PDF zeichenverschoben, entschlüsselt), S. 167 (4.2: Bezugsquerschnitt durch die 30°-Tangente);
%   Niemann/Winter 2003 S. 147 (21.7: Dauerbruch von der Zugseite an der 30°-Tangente, Bild 21.7/8); ISO 6336-3
%   S. 14 (6.1: chord between the 30° tangent contact points, Figures 3 and 4); Kassem 2022 S. 2 (Einleitung:
%   maximum at the 30° tangent section, notch effect); MA-0296 S. 8 (Leipersberger 42°-Tangente, 128 Brüche
%   POM/PA46, Hasl ohne erkennbaren Einfluss von Größe, Last, Temperatur; Abbildung 7).
% Absatz 3 (Schadensbilder): Hasl 2018 S. 93 (7.3: Anriss in der zugseitigen Fußrundung, keine Rastlinien,
%   Bild 37/38), S. 93 f. (Übergang zum Flankenversagen, kst-A mit Anschmelzungen), S. 100 (PA46 Übergang),
%   S. 135 (PEEK: Abscheren im Fußnutzkreis); MA-0296 S. 66 f. (6.2.1: 8 Nm vollständiger Bruch, 18 Nm Anriss,
%   <= 12 Nm vollständig, >= 14 Nm Anrisse, Risse links/rechts der Mitte, Stegrad, Tabelle 9), S. 76 (6.2.2
%   PEEK: Grübchen, Abscheren), S. 78 (6.2.3 PEEK-GF30: Grübchen, Verschleiß, freigelegte Faserenden).
% Bild (fig:quelle_fie23_zahnfussbruch): MA-0296 S. 66, Abbildung 54 (8 Nm, 39,4 Mio. LW) und Abbildung 55
%   (18 Nm, 787 000 LW).
% Absatz 4 (Grübchen, Verschleiß, Anschmelzen, Verformung): Illenberger 2022 S. 9 (2.2.2 Grübchen, 2.2.3
%   Verschleiß), S. 9 f. (Verdrehflankenspiel, Ölschmierung), S. 10 (2.2.4 Anschmelzungen, geschmiert möglich);
%   Fürstenberger 2013 S. 73 (6.x: Ausfall innerhalb weniger Minuten, Abb. 6.4), S. 125 (thermische
%   Leistungsgrenze, Anschmelzungen, Ausfall); VDI 2736 Blatt 2 S. 19 f. (3.2.5 Verformung, lambda_zul <= 0,07 m_n);
%   Hasl 2018 S. 16 (Winkler: Übergang plastische Verformung zu Zahnfußbruch, 10^3 < N_L <= 4·10^4).
% Absatz 5 (Versagensdefinitionen): VDI 2736 Blatt 4 S. 6 (3 und 3.1: Lastwechsel bis Bruch bzw. Grübchenmenge,
%   Schadensart nach Abschluss), S. 8 (3.4: 30 % Flächenanteil), S. 10 (5.1: Drehzahlabfall, zwei Messzellen);
%   Hasl 2018 S. 91 (7.2: Loss-of-drive, Setzeffekte, Drehmoment vor dem Ausfall); MA-0296 S. 65 (6.1:
%   Loss-of-drive, schleichendes Risswachstum, > 1 Nm über 100 000 LW, Abbildung 53).
% Absatz 6 (Einflüsse Temperatur, Dynamik, Schmierung, Werkstoff): VDI 2736 Blatt 2 S. 8 (3.2.1 Zahnfußtemperatur
%   setzt Zahnfußtragfähigkeit herab), S. 14 f. (Bild 3, Tabelle 5 sigma_FlimN über theta_Fuß); Hasl 2018 S. 91
%   (7.1: 80 °C Referenz), S. 134 (10: PA46 unter POM bei 80 °C, 120 °C moderat geringer; Hauptresonanz);
%   Illenberger 2022 S. 10 (Schmierung: Reibung, Wärmeabfuhr), S. 24 (2.5: abrasiver Verschleiß bei Ölschmierung
%   reduziert), S. 22 (Tavčar: faserverstärkt bei hohen Lasten, unverstärkt bei hohen Lastspielzahlen), S. 106
%   (9.2: Biegewechselfestigkeit, CF-PEEK Grübchen und Verschleiß, Hybridräder); Hasl 2018 S. 135 (Schmier- und
%   Kühlbedingungen); FVA 856 I S. 128 (Randschicht, siehe Abschnitt 2.2).
% Absatz 7 (Geometrie, Qualität): Klocke/Brecher 2017 S. 93 (3.1.2: Modul bestimmt Zahnfußquerschnitt), S. 31
%   (2.2: Eingriffswinkel, Krümmungsradien, Zahnfußdickensehnen), S. 95 (Unterschnitt, Bild 3.10; PDF-Seite =
%   Druckseite + 17, geprüft an 6.4.1 S. 604 und Inhaltsverzeichnis); Linke 2010 S. 121 (Bild 2.3/14 Zahnfußdicke
%   an der 30°-Tangente, Schneidrad gegen Wälzfräser); Hasl 2018 S. 134 (Zähnezahlvariation, kst-A/kst-B über
%   VDI 2736); Kassem 2022 S. 2 (notch effect); VDI 2736 Blatt 1 S. 30 (8.1: Teilungs-, Profilwinkel-,
%   Rundlaufabweichung, Zahnformabweichung).
% Bildwunsch:
% fig:quelle_fie23_zahnfussbruch: MA-0296.txt / 01_beispielarbeiten/MA-0296.pdf, Druckseite 66 (PDF-Seite 79),
%   Abbildung 54 "Kompletter Zahnfußbruch eines kst-E-Zahnrads aus POM bei niedriger Last 8 Nm" und Abbildung 55
%   "Zahnfußanriss eines kst-E Zahnrads aus POM bei hoher Last 18 Nm"; beide Fotos nebeneinander ausschneiden.
%   Datei kap2/quelle_fie23_zahnfussbruch.
% Optional (nicht eingebunden): Illenberger 2022 Bild 2-1 "Exemplarischer Zahnfußbruch an einem POM-Prüfrad
%   [Has19]", S. 8 (PDF-Seite 22), zeigt den Anriss an der 30°-Tangente; Fürstenberger 2013 Abb. 6.4
%   "Exemplarische Anschmelzungen der Prüfräder aus PA12G", S. 73 (PDF-Seite 87), entspricht Bild 2.6 der BA.
% Bibwunsch: Fuer13, NieWin03 in <SCRATCH>\bib_kap2_2_4.bib; Ill22, Has18, Fie23, Lin10, KloBre17, ISO6336-3,
%   VDI2736-1, VDI2736-2, VDI2736-4, Kas22 stehen in litbib.bib.
% Nomenklatur: keine neuen Zeichen; m_n (Normalmodul) steht in Nomenklatur.tex.
% Offen:
% BA 2.2.3 Zeile 1035 ([Wal89], Bild 2.5 Zahnradschäden nach Walton): nicht vorhanden, Bild durch MA-0296
%   Abbildung 54/55 ersetzt, Schadensliste auf VDI 2736 Blatt 1 Tabelle 4 umgestellt.
% Zeile 1047 bis 1052 ([Kra84] Balkenmodell, Schub und Druck vernachlässigt): auf Illenberger S. 8 umgestellt,
%   der dieselbe Aussage nach Krause bringt.
% Zeile 1053 bis 1056 ([Lin10] Grübchen, Öl in Haarrisse): auf Illenberger S. 9 umgestellt (Lin10 dort als
%   [Lin96] zitiert; eigene Linke-Seite zu Grübchen nicht geprüft).
% Zeile 1057 bis 1061 ([Dud61] Verschleiß, Querschnitt, Bruch): auf Illenberger S. 9 umgestellt.
% Zeile 1064 bis 1069 ([Für13] Trockenlauf, Anschmelzen; [Rös19] thermische Ermüdung mit der Frequenz,
%   Hysterese): Fürstenberger S. 1, 73 und 125 geprüft und zitiert; [Rös19] nicht vorhanden, Aussage auf
%   Illenberger S. 10 (viskoelastische Verluste) verkürzt.
% Zeile 1081 bis 1090 ([SAE21] Mehrbereichsöle, Kaltstart von Motorrädern): gestrichen, nicht zum Thema.
% Zeile 1093 bis 1103 ([ISO14] Makro- und Mikrogeometrie, Bild 2.7 Standardkorrekturen): gestrichen, Bild nicht
%   übernommen; Verzahnungsabweichungen über VDI 2736 Blatt 1 S. 30.
% Zeile 1104 bis 1128 ([Klo17, Nie03] Zähnezahl und Modul, Eingriffswinkel, Schrägungswinkel): auf
%   Klocke/Brecher S. 31, 93 und 95 verkürzt; Schrägungswinkel (Sprungüberdeckung, Axialkraft) gestrichen, weil
%   die Prüfverzahnung geradverzahnt ist; Nie03 als NieWin03 (2003, zweiter berichtigter Nachdruck) nur für die
%   Schadensarten und die 30°-Tangente zitiert.
% Zeile 1129 bis 1135 ([Rot98] komplementäre Zahndicken, Profilverschiebung): gestrichen, Roth 1998 nicht
%   vorhanden; der Profilverschiebungseinfluss bleibt ohne geprüfte Quelle offen (Klocke/Brecher S. 59 behandelt
%   nur Kegelräder).
% Zeile 1136 bis 1140 (Zielkonflikte NVH, Wirkungsgrad): gestrichen, ohne Quelle.
% Fürstenberger: das PDF ist der Entwurf "Stand 01.03.2013" mit Platzhaltern für die Einreichdaten; Seitenzahlen
%   nach diesem Entwurf.
